% ============================================================
%  SOAL + PEMBAHASAN  --  MATEMATIKA  --  SSU-ITB 2024
%  Tipe A (1-20) + Tipe B (21-39) + Soal Baru (40) | Paket 1
% ============================================================

\documentclass[11pt]{article}

\usepackage[utf8]{inputenc}
\usepackage[T1]{fontenc}
\usepackage[a4paper,margin=2.5cm,top=2.8cm]{geometry}
\usepackage{amsmath,amssymb}
\usepackage{graphicx}
\usepackage{enumitem}
\usepackage{multicol}
\usepackage{xcolor}
\usepackage[most]{tcolorbox}
\usepackage{fancyhdr}
\usepackage{booktabs}
\usepackage{icomma}
\usepackage{array}

% --- Warna Matematika ---
\definecolor{mapelcolor}{RGB}{20,110,60}
\definecolor{mapelbg}{RGB}{228,246,234}
\definecolor{pembcolor}{RGB}{100,50,130}
\definecolor{pembbg}{RGB}{245,238,252}
\definecolor{gold}{RGB}{210,160,30}

\pagestyle{fancy}\fancyhf{}
\lhead{\small SSU-ITB 2024 --- Pembahasan}
\rhead{\small Paket 1}
\cfoot{\small\thepage}
\renewcommand{\headrulewidth}{0.4pt}

\newtcolorbox{soal}[1]{
  breakable,
  colback=mapelbg, colframe=mapelcolor,
  coltitle=white, fonttitle=\bfseries\sffamily,
  title=Nomor #1,
  boxrule=0.8pt, arc=3pt,
  left=3mm, right=3mm, top=2mm, bottom=2mm}

\newtcolorbox{pembox}{
  breakable,
  colback=pembbg, colframe=pembcolor,
  coltitle=white, fonttitle=\bfseries\sffamily,
  title=Pembahasan,
  boxrule=0.8pt, arc=3pt,
  left=3mm, right=3mm, top=2mm, bottom=2mm}

\newcommand{\jawaban}[1]{%
  \noindent\colorbox{gold!30}{\parbox{\dimexpr\linewidth-2\fboxsep\relax}{%
    \textbf{Jawaban:} #1}}\par\smallskip}

\newcommand{\konsep}{\noindent\textit{Konsep:}\ \ignorespaces}

\newenvironment{pil}{%
  \begin{enumerate}[label=(\Alph*),leftmargin=*,itemsep=1pt,topsep=3pt]}%
  {\end{enumerate}}
\newenvironment{pilB}{%
  \par\smallskip\begin{multicols}{2}
  \begin{enumerate}[label=(\Alph*),leftmargin=*,itemsep=1pt,topsep=3pt]}%
  {\end{enumerate}\end{multicols}}

\newcommand{\gambar}[2][0.52\textwidth]{%
  \begin{center}\includegraphics[width=#1]{#2}\end{center}}

\linespread{1.05}

% ===================================================================
\begin{document}
\thispagestyle{empty}

\begin{center}
  \fbox{\parbox{0.6\textwidth}{\centering\bfseries\large
    LEMBAR SOAL DAN PEMBAHASAN}}
\end{center}
\vspace{0.5cm}

\begin{center}
  {\LARGE\bfseries MATEMATIKA}
\end{center}
\vspace{0.5cm}

\begin{center}
  \begin{tabular}{@{\hspace{2cm}}ll@{\hspace{0.4cm}:@{\hspace{0.4cm}}}l}
    Jenjang      & & SMA / Sederajat               \\[0.3em]
    Hari/Tanggal & & \underline{\hspace{5cm}}      \\[0.3em]
    Waktu        & & \underline{\hspace{5cm}}      \\[0.3em]
    Paket Soal   & & 1                             \\
  \end{tabular}
\end{center}

\vspace{0.6cm}
\noindent\rule{\linewidth}{0.6pt}
\vspace{0.3cm}

\noindent\textbf{\underline{PETUNJUK UMUM}}
\begin{enumerate}[leftmargin=*,itemsep=4pt,topsep=4pt]
  \item Anda \textbf{tidak} diperbolehkan menggunakan sumber-sumber eksternal,
        termasuk internet, \textit{large language model} seperti ChatGPT, Claude,
        LLaMA, Gemini, Meta AI, dan sebagainya. Anda sangat dianjurkan untuk
        mencoba mengerjakan sendiri terlebih dahulu karena evaluasi ini dibuat
        untuk \textbf{mengukur} tingkat kepahaman; nilai $<70$ mengindikasikan
        \textbf{tidak lulus}, dan jika sudah melewati \textit{deadline} maka
        dianggap tidak mengerjakan yang berarti \textbf{tidak lulus}.

  \item Dokumen ini merupakan \textbf{lembar soal sekaligus pembahasan}; gunakan
        pada tahap \textbf{Review}. Di tahap ini Anda diharapkan mengerjakan ulang
        dengan disertakan penjelasan/pembelaan mengapa revisi jawaban adalah pilihan
        tersebut, dan tetap mencantumkan penjelasan pada soal yang walaupun pada
        penilaian sudah benar.

  \item Anda dianjurkan untuk mencantumkan caranya walaupun tidak termasuk
        penilaian, mengingat sifatnya adalah sebagai bahan pembelajaran.
        Mohon bertanggung jawab atas pekerjaan Anda sendiri.

  \item Terdapat tiga jenis soal: \textbf{Pilihan Ganda Biasa} (5 opsi, pilih 1);
        \textbf{Pilihan Ganda Majemuk Kompleks} (tabel \textbf{Benar}/\textbf{Salah},
        tiap pernyataan dinilai mandiri); dan \textbf{Isian Singkat}
        (tulis nilai/ekspresi).

  \item Soal berjumlah \textbf{40 butir}: nomor 1--20 dari SSU-ITB~2024 Tipe~A,
        nomor 21--39 dari Tipe~B, dan nomor 40 adalah soal tambahan.

  \item Isilah detail informasi berikut.\\[0.3em]
        \textbf{Nama}\hspace{0.45cm};\enspace\underline{\hspace{8cm}}
\end{enumerate}

\vspace{0.6cm}
\noindent\textit{Dengan menandatangani kolom di bawah ini, saya menyatakan telah membaca,
memahami, dan menyetujui seluruh peraturan yang tercantum di atas.}

\vspace{0.6cm}
\noindent\fbox{\parbox{5cm}{\vspace{2.2cm}\centering\small Tanda Tangan Peserta\vspace{0.3cm}}}

\clearpage

% ===================================================================
%  TIPE A  (nomor 1--20)
% ===================================================================

\begin{soal}{1}
Jika titik $P(a,b)$ berada pada kurva $y=x^3+5x-7$ di kuadran pertama
sehingga gradien garis singgung kurva tersebut di $P$ adalah 17,
maka nilai $a+b$ adalah \ldots
\begin{pilB}
  \item $9$ \item $11$ \item $13$ \item $15$ \item $17$
\end{pilB}
\end{soal}
\begin{pembox}
\jawaban{(C) $a+b=13$}
\konsep Gradien garis singgung kurva $y=f(x)$ di titik $P$ adalah $f'(x_P)$.

Turunan kurva: $y'=3x^2+5$. Di $P(a,b)$:
\[ 3a^2+5=17 \;\Rightarrow\; a^2=4 \;\Rightarrow\; a=2 \text{ (kuadran I, }a>0\text{)}. \]
Nilai $b$: $b=2^3+5(2)-7=8+10-7=11$. Jadi $a+b=2+11=13$.
\end{pembox}

\begin{soal}{2}
Perhatikan gambar berikut.
\gambar[0.46\textwidth]{images/1/1-17.png}
\begin{enumerate}[label=(\alph*),leftmargin=*,itemsep=6pt,topsep=4pt]
  \item Jika $L(a)$ menyatakan luas persegi panjang di atas,
        maka $L(8)$ adalah \ldots
  \item Jika $c$ memenuhi $L'(c)=0$, maka nilai $2\sqrt{2}\,c$ adalah \ldots
  \item Luas persegi panjang maksimum yang mungkin termuat pada
        daerah setengah lingkaran di atas adalah \ldots
\end{enumerate}
\end{soal}
\begin{pembox}
\jawaban{(a) $L(8)=96$\quad (b) $2\sqrt{2}\,c=20$\quad (c) $L_{\max}=100$}
\konsep Setengah lingkaran $x^2+y^2=100$; persegi panjang dengan sudut kanan atas di $(a,\sqrt{100-a^2})$, lebar $2a$, tinggi $\sqrt{100-a^2}$.
\[ L(a)=2a\sqrt{100-a^2}. \]
\textbf{(a)} $L(8)=16\sqrt{36}=96$.

\textbf{(b)} $L'(a)=\dfrac{200-4a^2}{\sqrt{100-a^2}}=0\Rightarrow c=5\sqrt{2}$. Maka $2\sqrt{2}\cdot5\sqrt{2}=20$.

\textbf{(c)} $L(5\sqrt{2})=2(5\sqrt{2})(5\sqrt{2})=100$.
\end{pembox}

\begin{soal}{3}
Dua bola diambil sekaligus dari suatu kotak yang berisi 5 bola biru
dan 6 bola kuning. Jika peluang terambil paling banyak 1 bola kuning
adalah $\dfrac{A}{11\times10}$, maka nilai $A$ adalah \ldots
\begin{pilB}
  \item $40$ \item $60$ \item $80$ \item $90$ \item $100$
\end{pilB}
\end{soal}
\begin{pembox}
\jawaban{(C) $A=80$}
\konsep $P(\text{paling banyak 1 kuning})=P(0\text{ kuning})+P(1\text{ kuning})$.
\[ P=\frac{\binom{5}{2}+\binom{6}{1}\binom{5}{1}}{\binom{11}{2}}=\frac{10+30}{55}=\frac{8}{11}=\frac{80}{110}=\frac{A}{11\times10}. \]
Jadi $A=80$.
\end{pembox}

\begin{soal}{4}
Diberikan empat pernyataan tentang hiperbola $y^2-x^2=1$.
Tentukan kebenaran setiap pernyataan berikut!

\par\smallskip
{\renewcommand{\arraystretch}{1.4}%
\begin{tabular}{@{}p{0.76\linewidth}cc@{}}
\hline
\textbf{Pernyataan} & \textbf{Benar} & \textbf{Salah}\\
\hline
(1)~Hiperbola $y^2-x^2=1$ memotong sumbu $y$ di dua titik.
  & $\bigcirc$ & $\bigcirc$\\
(2)~Hiperbola $y^2-x^2=1$ merupakan grafik suatu fungsi.
  & $\bigcirc$ & $\bigcirc$\\
(3)~Hiperbola $y^2-x^2=1$ tidak memotong sumbu $x$.
  & $\bigcirc$ & $\bigcirc$\\
(4)~Hiperbola $y^2-x^2=1$ tidak memiliki asimtot miring.
  & $\bigcirc$ & $\bigcirc$\\
\hline
\end{tabular}}
\end{soal}
\begin{pembox}
\jawaban{(1)~\textbf{Benar}\quad(2)~\textbf{Salah}\quad(3)~\textbf{Benar}\quad(4)~\textbf{Salah}}
\begin{enumerate}[leftmargin=*,topsep=3pt,itemsep=3pt]
\item \textbf{Benar.} Saat $x=0$: $y^2=1\Rightarrow y=\pm1$. Dua titik potong.
\item \textbf{Salah.} Untuk setiap $x$, $y=\pm\sqrt{x^2+1}$; tidak memenuhi syarat fungsi.
\item \textbf{Benar.} Saat $y=0$: $-x^2=1$, tidak ada solusi real.
\item \textbf{Salah.} Hiperbola ini memiliki asimtot miring $y=x$ dan $y=-x$.
\end{enumerate}
\end{pembox}

\begin{soal}{5}
Jika $\cos(2x)=-\dfrac{7}{25}$ dengan $0<x<\pi$,
maka $25\sin(x)$ adalah \ldots
\begin{pilB}
  \item $12$ \item $15$ \item $16$ \item $20$ \item $24$
\end{pilB}
\end{soal}
\begin{pembox}
\jawaban{(D) $25\sin x=20$}
\konsep $\cos 2x=1-2\sin^2 x$.
\[ 1-2\sin^2 x=-\tfrac{7}{25}\;\Rightarrow\;\sin^2 x=\tfrac{16}{25}\;\Rightarrow\;\sin x=\tfrac{4}{5}\;(0<x<\pi). \]
Jadi $25\sin x=25\cdot\tfrac{4}{5}=20$.
\end{pembox}

\begin{soal}{6}
Diberikan titik $A(1,-3,3)$ dan titik $B(-5,-2,3)$. Jika
$P(p_1,p_2,p_3)$ adalah suatu titik sehingga
$\overrightarrow{BP}=-4\overrightarrow{PA}$,
maka nilai $-3p_2$ adalah \ldots
\begin{pilB}
  \item $-15$ \item $-10$ \item $5$ \item $10$ \item $15$
\end{pilB}
\end{soal}
\begin{pembox}
\jawaban{(D) $-3p_2=10$}
$P-B=-4(A-P)\;\Rightarrow\;3P=4A-B$.
\[ P=\frac{4(1,-3,3)-(-5,-2,3)}{3}=\frac{(9,-10,9)}{3}=\!\left(3,-\tfrac{10}{3},3\right). \]
Jadi $-3p_2=-3\!\left(-\tfrac{10}{3}\right)=10$.
\end{pembox}

\begin{soal}{7}
Diketahui suku ke-10 dan suku ke-20 suatu barisan aritmetika,
berturut-turut, adalah 84 dan 174.
Tentukan kebenaran setiap pernyataan berikut!

\par\smallskip
{\renewcommand{\arraystretch}{1.4}%
\begin{tabular}{@{}p{0.76\linewidth}cc@{}}
\hline
\textbf{Pernyataan} & \textbf{Benar} & \textbf{Salah}\\
\hline
(1)~Beda $(d)$ barisan tersebut adalah 18.
  & $\bigcirc$ & $\bigcirc$\\
(2)~Suku pertama barisan tersebut adalah 3.
  & $\bigcirc$ & $\bigcirc$\\
(3)~Suku ke-15 barisan tersebut adalah 129.
  & $\bigcirc$ & $\bigcirc$\\
(4)~Jumlah 15 suku pertama $S_{15}=990$.
  & $\bigcirc$ & $\bigcirc$\\
\hline
\end{tabular}}
\end{soal}
\begin{pembox}
\jawaban{(1)~\textbf{Salah}\quad(2)~\textbf{Benar}\quad(3)~\textbf{Benar}\quad(4)~\textbf{Benar}}
$U_{20}-U_{10}=10d=90\Rightarrow d=9$ (bukan 18).
$a_1=84-9\cdot9=3$.
$U_{15}=3+14(9)=129$.
$S_{15}=\tfrac{15}{2}(2\cdot3+14\cdot9)=\tfrac{15}{2}(132)=990$.
\end{pembox}

\begin{soal}{8}
Perhatikan fungsi $f(x)=\sqrt{8-|3x-7|}$.
Tentukan kebenaran setiap pernyataan berikut!

\par\smallskip
{\renewcommand{\arraystretch}{1.4}%
\begin{tabular}{@{}p{0.76\linewidth}cc@{}}
\hline
\textbf{Pernyataan} & \textbf{Benar} & \textbf{Salah}\\
\hline
(1)~$x=0$ termasuk dalam domain fungsi $f$.
  & $\bigcirc$ & $\bigcirc$\\
(2)~$x=5$ termasuk dalam domain fungsi $f$.
  & $\bigcirc$ & $\bigcirc$\\
(3)~$x=-1$ termasuk dalam domain fungsi $f$.
  & $\bigcirc$ & $\bigcirc$\\
(4)~Domain fungsi $f$ adalah $\left[-\tfrac{1}{3},\,5\right]$.
  & $\bigcirc$ & $\bigcirc$\\
\hline
\end{tabular}}
\end{soal}
\begin{pembox}
\jawaban{(1)~\textbf{Benar}\quad(2)~\textbf{Benar}\quad(3)~\textbf{Salah}\quad(4)~\textbf{Benar}}
$|3x-7|\leq8\;\Rightarrow\;-1\leq3x\leq15\;\Rightarrow\;x\in\!\left[-\tfrac{1}{3},5\right]$.
(1) $0\in[-\tfrac{1}{3},5]$\;\checkmark\quad
(2) $5\in[-\tfrac{1}{3},5]$\;\checkmark\quad
(3) $-1<-\tfrac{1}{3}$, di luar domain.\quad
(4) Benar.
\end{pembox}

\begin{soal}{9}
Diberikan $f(x)={}^{13}\!\log(x)$.
Tentukan kebenaran setiap pernyataan berikut!

\par\smallskip
{\renewcommand{\arraystretch}{1.4}%
\begin{tabular}{@{}p{0.76\linewidth}cc@{}}
\hline
\textbf{Pernyataan} & \textbf{Benar} & \textbf{Salah}\\
\hline
(1)~$f(5x)+f\!\left(\dfrac{5}{x}\right)={}^{13}\!\log 25$.
  & $\bigcirc$ & $\bigcirc$\\[6pt]
(2)~$f(5x)+f\!\left(\dfrac{5}{x}\right)=2\cdot{}^{13}\!\log 5$.
  & $\bigcirc$ & $\bigcirc$\\[6pt]
(3)~$f(5x)+f\!\left(\dfrac{5}{x}\right)={}^{13}\!\log 10$.
  & $\bigcirc$ & $\bigcirc$\\[6pt]
(4)~$f(169)=2$.
  & $\bigcirc$ & $\bigcirc$\\
\hline
\end{tabular}}
\end{soal}
\begin{pembox}
\jawaban{(1)~\textbf{Benar}\quad(2)~\textbf{Benar}\quad(3)~\textbf{Salah}\quad(4)~\textbf{Benar}}
$f(5x)+f\!\left(\tfrac{5}{x}\right)=\log_{13}(5x)+\log_{13}\!\left(\tfrac{5}{x}\right)=\log_{13}25=2\log_{13}5$.
Bukan $\log_{13}10$. Dan $f(169)=\log_{13}(13^2)=2$. \checkmark
\end{pembox}

\begin{soal}{10}
Data 90 dari 100 mahasiswa baru suatu perguruan tinggi disajikan
pada tabel berikut.
\begin{center}
\begin{tabular}{l|ccc}
\hline
Jalur   & Kota A & Kota B & Selain A dan B \\
\hline
SNMPTN  & 10 & 12 &  8 \\
UTBK    &  9 &  4 & 17 \\
Mandiri & 11 & 14 &  5 \\
\hline
\end{tabular}
\end{center}
Adapun 10 mahasiswa sisanya bukan dari Kota A dan bukan jalur SNMPTN.
Misalkan $X$ adalah mahasiswa jalur UTBK yang dipilih acak dari 100
mahasiswa. Tentukan kebenaran setiap pernyataan berikut!

\par\smallskip
{\renewcommand{\arraystretch}{1.4}%
\begin{tabular}{@{}p{0.76\linewidth}cc@{}}
\hline
\textbf{Pernyataan} & \textbf{Benar} & \textbf{Salah}\\
\hline
(1)~Dari 90 data, jumlah mahasiswa UTBK asal Kota B adalah 4.
  & $\bigcirc$ & $\bigcirc$\\
(2)~Dari 90 data, setiap jalur memiliki tepat 30 mahasiswa.
  & $\bigcirc$ & $\bigcirc$\\
(3)~Peluang $X$ berasal dari Kota B dapat dipastikan $<\tfrac{1}{12}$.
  & $\bigcirc$ & $\bigcirc$\\
(4)~Dari 90 data, jumlah dari Kota A sama dengan jumlah dari Kota B.
  & $\bigcirc$ & $\bigcirc$\\
\hline
\end{tabular}}
\end{soal}
\begin{pembox}
\jawaban{(1)~\textbf{Benar}\quad(2)~\textbf{Benar}\quad(3)~\textbf{Salah}\quad(4)~\textbf{Benar}}
(1) Dari tabel: UTBK Kota B $=4$.\;(2) SNMPTN$=30$, UTBK$=30$, Mandiri$=30$.
(3) 10 mahasiswa sisa bisa masuk UTBK Kota B, sehingga jumlah total UTBK belum pasti; peluang tidak dapat dipastikan $<\tfrac{1}{12}$.
(4) Kota A: $10{+}9{+}11=30$; Kota B: $12{+}4{+}14=30$. Sama.
\end{pembox}

\begin{soal}{11}
Diketahui suatu parabola grafiknya melalui $(-4,0)$ dan $(4,0)$,
serta titik maksimumnya terletak pada garis $y=23x+48$.
Tentukan kebenaran setiap pernyataan berikut!

\par\smallskip
{\renewcommand{\arraystretch}{1.4}%
\begin{tabular}{@{}p{0.76\linewidth}cc@{}}
\hline
\textbf{Pernyataan} & \textbf{Benar} & \textbf{Salah}\\
\hline
(1)~Koefisien $x^2$ pada persamaan parabola bernilai negatif.
  & $\bigcirc$ & $\bigcirc$\\
(2)~Persamaan parabola tersebut adalah $y=-3x^2+48$.
  & $\bigcirc$ & $\bigcirc$\\
(3)~Titik puncak parabola berada di $(0,24)$.
  & $\bigcirc$ & $\bigcirc$\\
(4)~Parabola memotong sumbu $x$ di $(-4,0)$ dan $(4,0)$.
  & $\bigcirc$ & $\bigcirc$\\
\hline
\end{tabular}}
\end{soal}
\begin{pembox}
\jawaban{(1)~\textbf{Benar}\quad(2)~\textbf{Benar}\quad(3)~\textbf{Salah}\quad(4)~\textbf{Benar}}
Bentuk: $y=k(x^2-16)$. Sumbu simetri $x=0$; titik puncak saat $x=0$ pada garis: $y=48$. Maka $-16k=48\Rightarrow k=-3$. Jadi $y=-3x^2+48$, puncak di $(0,48)$ bukan $(0,24)$.
\end{pembox}

\begin{soal}{12}
Tentukan nilai
\[
  \lim_{x\to 3}\frac{\sin(21-7x)}{x-3}.
\]
\begin{pilB}
  \item $-21$ \item $-7$ \item $0$ \item $7$ \item $21$
\end{pilB}
\end{soal}
\begin{pembox}
\jawaban{(B) $-7$}
\konsep $\lim_{u\to 0}\dfrac{\sin u}{u}=1$.

$21-7x=-7(x-3)\to0$ saat $x\to3$. Maka
\[ \lim_{x\to 3}\frac{\sin(-7(x-3))}{x-3}=\lim\frac{-7(x-3)}{x-3}=-7. \]
\end{pembox}

\begin{soal}{13}
Jika
\[
  \int_{-7}^{-3}f(x)\,dx=12,
\]
maka nilai
\[
  \int_{3}^{5}f(3-2x)\,dx
\]
adalah \ldots
\begin{pilB}
  \item $2$ \item $4$ \item $6$ \item $8$ \item $12$
\end{pilB}
\end{soal}
\begin{pembox}
\jawaban{(C) $6$}
\konsep Substitusi $u=3-2x$, $du=-2\,dx$.

Batas: $x=3\Rightarrow u=-3$;\; $x=5\Rightarrow u=-7$.
\[ \int_3^5 f(3-2x)\,dx=\int_{-3}^{-7}f(u)\!\left(-\tfrac{1}{2}\right)du=\tfrac{1}{2}\int_{-7}^{-3}f(u)\,du=\tfrac{1}{2}(12)=6. \]
\end{pembox}

\begin{soal}{14}
Perhatikan grafik fungsi $f$ berikut.
\gambar[0.5\textwidth]{images/1/1-18.png}
Diketahui $a=6$. Tentukan kebenaran setiap pernyataan berikut!

\par\smallskip
{\renewcommand{\arraystretch}{1.4}%
\begin{tabular}{@{}p{0.76\linewidth}cc@{}}
\hline
\textbf{Pernyataan} & \textbf{Benar} & \textbf{Salah}\\
\hline
(1)~$f(6)=36$.
  & $\bigcirc$ & $\bigcirc$\\
(2)~$f'(6)=0$, karena $a=6$ adalah titik maksimum fungsi $f$.
  & $\bigcirc$ & $\bigcirc$\\
(3)~$\displaystyle\lim_{x\to 6}\dfrac{f(x)-6x}{x-6}=f'(6)-6$.
  & $\bigcirc$ & $\bigcirc$\\
(4)~$\displaystyle\lim_{x\to 6}\dfrac{f(x)-6x}{x-6}=6$.
  & $\bigcirc$ & $\bigcirc$\\
\hline
\end{tabular}}
\end{soal}
\begin{pembox}
\jawaban{(1)~\textbf{Benar}\quad(2)~\textbf{Benar}\quad(3)~\textbf{Benar}\quad(4)~\textbf{Salah}}
Dari grafik, puncak di $(a,a^2)=(6,36)$, sehingga $f(6)=36=6\cdot6$. Karena $a=6$ adalah titik maksimum, $f'(6)=0$.
\[ \lim_{x\to6}\frac{f(x)-6x}{x-6}=\lim_{x\to6}\frac{f(x)-f(6)}{x-6}-\lim_{x\to6}\frac{6x-6a}{x-6}=f'(6)-6=0-6=-6\neq6. \]
\end{pembox}

\begin{soal}{15}
Diketahui $x^3-4x^2+9x+a=(x-1)\,g(x)$ dengan $a$ konstanta real.
Tentukan kebenaran setiap pernyataan berikut!

\par\smallskip
{\renewcommand{\arraystretch}{1.4}%
\begin{tabular}{@{}p{0.76\linewidth}cc@{}}
\hline
\textbf{Pernyataan} & \textbf{Benar} & \textbf{Salah}\\
\hline
(1)~Nilai $a=-6$.
  & $\bigcirc$ & $\bigcirc$\\
(2)~Nilai $g(1)=4$.
  & $\bigcirc$ & $\bigcirc$\\
(3)~Polinom $x^3-4x^2+9x-6$ habis dibagi $(x-1)$.
  & $\bigcirc$ & $\bigcirc$\\
(4)~Nilai $g(0)=-6$.
  & $\bigcirc$ & $\bigcirc$\\
\hline
\end{tabular}}
\end{soal}
\begin{pembox}
\jawaban{(1)~\textbf{Benar}\quad(2)~\textbf{Benar}\quad(3)~\textbf{Benar}\quad(4)~\textbf{Salah}}
$P(1)=0$: $1-4+9+a=0\Rightarrow a=-6$. $g(1)=P'(1)=3-8+9=4$.
Bagi $P(x)$ oleh $(x-1)$: $g(x)=x^2-3x+6$, sehingga $g(0)=6$ (bukan $-6$).
\end{pembox}

\begin{soal}{16}
Diketahui fungsi $y=6\cos\!\left(\dfrac{\pi t}{3}\right)+9$.
Tentukan kebenaran setiap pernyataan berikut!

\par\smallskip
{\renewcommand{\arraystretch}{1.4}%
\begin{tabular}{@{}p{0.76\linewidth}cc@{}}
\hline
\textbf{Pernyataan} & \textbf{Benar} & \textbf{Salah}\\
\hline
(1)~Periode fungsi tersebut adalah 6.
  & $\bigcirc$ & $\bigcirc$\\
(2)~Amplitudo fungsi tersebut adalah 9.
  & $\bigcirc$ & $\bigcirc$\\
(3)~Nilai minimum fungsi tersebut adalah 3.
  & $\bigcirc$ & $\bigcirc$\\
(4)~Jika $P$ menyatakan periode dan $A$ nilai minimum, maka $10P+A=63$.
  & $\bigcirc$ & $\bigcirc$\\
\hline
\end{tabular}}
\end{soal}
\begin{pembox}
\jawaban{(1)~\textbf{Benar}\quad(2)~\textbf{Salah}\quad(3)~\textbf{Benar}\quad(4)~\textbf{Benar}}
$\omega=\tfrac{\pi}{3}$. Periode $P=\tfrac{2\pi}{\pi/3}=6$. Amplitudo $=6$ (bukan 9; angka 9 adalah geseran vertikal). Minimum $=6(-1)+9=3$. Maka $10(6)+3=63$.
\end{pembox}

\begin{soal}{17}
Jika $f(x)=(a+1)x^3+x^2+12$ dan $f'(1)=32$,
maka nilai $a$ adalah \ldots
\begin{pilB}
  \item $5$ \item $7$ \item $8$ \item $9$ \item $11$
\end{pilB}
\end{soal}
\begin{pembox}
\jawaban{(D) $a=9$}
$f'(x)=3(a+1)x^2+2x$. Substitusi $x=1$:
\[ 3(a+1)+2=32\;\Rightarrow\;3a+5=32\;\Rightarrow\;a=9. \]
\end{pembox}

\begin{soal}{18}
Diketahui $f(x)=ax^2+b$,\;
$\displaystyle\int_0^1 f(x)\,dx=6$,\;
dan $\displaystyle\int_1^2 f(x)\,dx=12$.
Tentukan kebenaran setiap pernyataan berikut!

\par\smallskip
{\renewcommand{\arraystretch}{1.4}%
\begin{tabular}{@{}p{0.76\linewidth}cc@{}}
\hline
\textbf{Pernyataan} & \textbf{Benar} & \textbf{Salah}\\
\hline
(1)~Nilai $a=3$.
  & $\bigcirc$ & $\bigcirc$\\
(2)~Nilai $b=5$.
  & $\bigcirc$ & $\bigcirc$\\
(3)~$f(1)=8$.
  & $\bigcirc$ & $\bigcirc$\\
(4)~$\displaystyle\int_0^2 f(x)\,dx=24$.
  & $\bigcirc$ & $\bigcirc$\\
\hline
\end{tabular}}
\end{soal}
\begin{pembox}
\jawaban{(1)~\textbf{Benar}\quad(2)~\textbf{Benar}\quad(3)~\textbf{Benar}\quad(4)~\textbf{Salah}}
$\int_0^1=\tfrac{a}{3}+b=6$;\; $\int_1^2=\tfrac{7a}{3}+b=12$. Kurangkan: $2a=6\Rightarrow a=3$, $b=5$.
$f(1)=3+5=8$.\; $\int_0^2=6+12=18\neq24$.
\end{pembox}

\begin{soal}{19}
Diberikan transformasi $T$ dengan matriks penyajian $M$ ($2\times2$),
yaitu $T(X)=MX$ untuk tiap $X=\begin{bmatrix}x\\y\end{bmatrix}$.
Diketahui
\[
  T\!\left(\begin{bmatrix}3\\-2\end{bmatrix}\right)=\begin{bmatrix}1\\1\end{bmatrix}
  \quad\text{dan}\quad
  T\!\left(\begin{bmatrix}-1\\1\end{bmatrix}\right)=\begin{bmatrix}4\\2\end{bmatrix}.
\]
Jika $T(X)=\begin{bmatrix}9\\5\end{bmatrix}$, tentukan kebenaran
pernyataan-pernyataan berikut tentang $X$!

\par\smallskip
{\renewcommand{\arraystretch}{1.4}%
\begin{tabular}{@{}p{0.76\linewidth}cc@{}}
\hline
\textbf{Pernyataan} & \textbf{Benar} & \textbf{Salah}\\
\hline
(1)~Komponen pertama dari $X$ adalah 1.
  & $\bigcirc$ & $\bigcirc$\\
(2)~Komponen kedua dari $X$ adalah 1.
  & $\bigcirc$ & $\bigcirc$\\
(3)~$X=1\cdot\begin{bmatrix}3\\-2\end{bmatrix}+2\cdot\begin{bmatrix}-1\\1\end{bmatrix}$.
  & $\bigcirc$ & $\bigcirc$\\
(4)~Norma vektor $X$ adalah $|X|=\sqrt{2}$.
  & $\bigcirc$ & $\bigcirc$\\
\hline
\end{tabular}}
\end{soal}
\begin{pembox}
\jawaban{(1)~\textbf{Benar}\quad(2)~\textbf{Salah}\quad(3)~\textbf{Benar}\quad(4)~\textbf{Salah}}
\konsep Linearitas: $T(\alpha\mathbf{u}+\beta\mathbf{v})=\alpha T(\mathbf{u})+\beta T(\mathbf{v})$.

Misalkan $X=\alpha\begin{bmatrix}3\\-2\end{bmatrix}+\beta\begin{bmatrix}-1\\1\end{bmatrix}$. Maka $T(X)=\alpha\begin{bmatrix}1\\1\end{bmatrix}+\beta\begin{bmatrix}4\\2\end{bmatrix}=\begin{bmatrix}9\\5\end{bmatrix}$.
Sistem: $\alpha+4\beta=9$, $\alpha+2\beta=5$ $\Rightarrow$ $\beta=2,\alpha=1$.
$X=\begin{bmatrix}3\\-2\end{bmatrix}+2\begin{bmatrix}-1\\1\end{bmatrix}=\begin{bmatrix}1\\0\end{bmatrix}$.
Komponen kedua $=0$ (bukan 1). $|X|=1$ (bukan $\sqrt{2}$).
\end{pembox}

\begin{soal}{20}
Diberikan limas persegi $T.ABCD$ seperti pada gambar berikut.
\gambar[0.52\textwidth]{images/1/1-19.png}
Diketahui $|AB|=2\sqrt{2}$ dan tinggi limas adalah 4. Titik $P$
adalah titik tengah $\overline{AT}$ dan titik $Q$ pada $\overline{CT}$
dengan $|CT|=4|CQ|$. Tentukan kebenaran setiap pernyataan berikut!

\par\smallskip
{\renewcommand{\arraystretch}{1.4}%
\begin{tabular}{@{}p{0.76\linewidth}cc@{}}
\hline
\textbf{Pernyataan} & \textbf{Benar} & \textbf{Salah}\\
\hline
(1)~Titik $P$ berada tepat di tengah ruas garis $\overline{AT}$.
  & $\bigcirc$ & $\bigcirc$\\
(2)~$|CQ|=\tfrac{1}{4}|CT|$, sehingga $Q$ seperempat perjalanan dari $C$ ke $T$.
  & $\bigcirc$ & $\bigcirc$\\
(3)~Garis $\overline{PQ}$ sejajar dengan bidang alas $ABCD$.
  & $\bigcirc$ & $\bigcirc$\\
(4)~Tangen sudut antara $\overline{AC}$ dan $\overline{PQ}$ adalah $\tfrac{2}{5}$.
  & $\bigcirc$ & $\bigcirc$\\
\hline
\end{tabular}}
\end{soal}
\begin{pembox}
\jawaban{(1)~\textbf{Benar}\quad(2)~\textbf{Benar}\quad(3)~\textbf{Salah}\quad(4)~\textbf{Benar}}
\konsep Koordinat: $A=(-\sqrt{2},-\sqrt{2},0)$, $C=(\sqrt{2},\sqrt{2},0)$, $T=(0,0,4)$.

$P=\frac{A+T}{2}=\!\left(-\frac{\sqrt{2}}{2},-\frac{\sqrt{2}}{2},2\right)$;\;
$Q=C+\frac{1}{4}(T-C)=\!\left(\frac{3\sqrt{2}}{4},\frac{3\sqrt{2}}{4},1\right)$.

$\overrightarrow{PQ}=\!\left(\frac{5\sqrt{2}}{4},\frac{5\sqrt{2}}{4},-1\right)$ --- komponen $z=-1\neq0$, sehingga $PQ$ tidak sejajar bidang alas.

$\tan\theta=\dfrac{|\vec{AC}\times\vec{PQ}|}{|\vec{AC}\cdot\vec{PQ}|}=\dfrac{2}{5}$.
\end{pembox}

% ===================================================================
%  TIPE B  (nomor 21--39)
% ===================================================================

\begin{soal}{21}
Diberikan garis $y=16-4x$ dan kurva $y=16-x^2$. Titik potong keduanya
adalah $(0,a)$ dan $(b,c)$, serta $L$ adalah luas daerah tertutup
yang dibatasi kurva tersebut.
Tentukan kebenaran setiap pernyataan berikut!

\par\smallskip
{\renewcommand{\arraystretch}{1.4}%
\begin{tabular}{@{}p{0.76\linewidth}cc@{}}
\hline
\textbf{Pernyataan} & \textbf{Benar} & \textbf{Salah}\\
\hline
(1)~Koordinat $y$ titik potong pertama (di $x=0$) adalah $a=16$.
  & $\bigcirc$ & $\bigcirc$\\
(2)~Koordinat $x$ titik potong kedua adalah $b=4$.
  & $\bigcirc$ & $\bigcirc$\\
(3)~Nilai $6L=48$.
  & $\bigcirc$ & $\bigcirc$\\
(4)~Nilai $a+b=20$.
  & $\bigcirc$ & $\bigcirc$\\
\hline
\end{tabular}}
\end{soal}
\begin{pembox}
\jawaban{(1)~\textbf{Benar}\quad(2)~\textbf{Benar}\quad(3)~\textbf{Salah}\quad(4)~\textbf{Benar}}
$16-4x=16-x^2\Rightarrow x(x-4)=0$. Titik: $(0,16)$ dan $(4,0)$. $a=16,b=4,a+b=20$.
$L=\int_0^4(4x-x^2)dx=\left[2x^2-\tfrac{x^3}{3}\right]_0^4=32-\tfrac{64}{3}=\tfrac{32}{3}$.
$6L=64$ (bukan 48).
\end{pembox}

\begin{soal}{22}
Diberikan balok $ABCD.EFGH$ dengan $|AB|=3$ dan $|BC|=|AE|=7$.
\gambar[0.5\textwidth]{images/1/1-20.png}
Misalkan $P$ adalah titik pada $BC$ dan $Q$ pada $CG$ sehingga
$|BP|=|GQ|=1$. Kosinus sudut antara garis $\overline{AP}$ dan
garis $\overline{HQ}$ adalah \ldots
\end{soal}
\begin{pembox}
\jawaban{$\cos\theta=\dfrac{9}{10}$}
Koordinat: $A=(0,0,0)$, $B=(3,0,0)$, $C=(3,7,0)$, $H=(0,7,7)$, $G=(3,7,7)$.
$P=(3,1,0)$ ($|BP|=1$), $Q=(3,7,6)$ ($|GQ|=1$).
$\vec{AP}=(3,1,0)$, $\vec{HQ}=(3,0,-1)$.
$\cos\theta=\dfrac{9}{\sqrt{10}\cdot\sqrt{10}}=\dfrac{9}{10}$.
\end{pembox}

\begin{soal}{23}
Misalkan $s$ adalah simpangan baku dari data terurut enam bilangan:
$4,\ 6-a,\ 6,\ 8,\ 8+a,\ 10$.
Tentukan kebenaran setiap pernyataan berikut!

\par\smallskip
{\renewcommand{\arraystretch}{1.4}%
\begin{tabular}{@{}p{0.76\linewidth}cc@{}}
\hline
\textbf{Pernyataan} & \textbf{Benar} & \textbf{Salah}\\
\hline
(1)~Rata-rata data tersebut adalah 7 untuk setiap nilai $a$.
  & $\bigcirc$ & $\bigcirc$\\
(2)~Nilai varians $s^2$ merupakan fungsi tidak turun terhadap $a\geq0$.
  & $\bigcirc$ & $\bigcirc$\\
(3)~Untuk $a=0$, nilai $s^2<6$.
  & $\bigcirc$ & $\bigcirc$\\
(4)~Untuk setiap $a$ yang memenuhi keterurutan data, selalu berlaku $s^2\leq6$.
  & $\bigcirc$ & $\bigcirc$\\
\hline
\end{tabular}}
\end{soal}
\begin{pembox}
\jawaban{(1)~\textbf{Benar}\quad(2)~\textbf{Benar}\quad(3)~\textbf{Benar}\quad(4)~\textbf{Salah}}
$\bar{x}=42/6=7$ (tetap untuk semua $a$). $s^2=\dfrac{\sum(x_i-7)^2}{6}=\dfrac{a^2+2a+11}{3}$.
$\dfrac{ds^2}{da}=\dfrac{2a+2}{3}>0$ (tidak turun). Untuk $a=0$: $s^2=\tfrac{11}{3}<6$.
Keterurutan mensyaratkan $0\leq a\leq2$. Saat $a=2$: $s^2=\tfrac{19}{3}>6$. Jadi (4) Salah.
\end{pembox}

\begin{soal}{24}
Jika $r_{1,2}=6\pm\sqrt{2}$ adalah akar-akar persamaan kuadrat
$x^2-Px+Q=0$, maka $P+Q$ adalah \ldots
\begin{pilB}
  \item $34$ \item $38$ \item $42$ \item $46$ \item $50$
\end{pilB}
\end{soal}
\begin{pembox}
\jawaban{(D) $P+Q=46$}
\konsep Vieta: $P=r_1+r_2$, $Q=r_1 r_2$.

$P=(6+\sqrt{2})+(6-\sqrt{2})=12$.\; $Q=(6+\sqrt{2})(6-\sqrt{2})=34$.
$P+Q=46$.
\end{pembox}

\begin{soal}{25}
Jika $y=x^2-16x+61$ dengan $x>8$, maka $x=C+\sqrt{y+D}$ dengan
$C$ dan $D$ merupakan dua konstanta bilangan real. Nilai $10C+D$
adalah \ldots
\begin{pilB}
  \item $75$ \item $80$ \item $83$ \item $86$ \item $90$
\end{pilB}
\end{soal}
\begin{pembox}
\jawaban{(C) $10C+D=83$}
$y=(x-8)^2-3\;\Rightarrow\;y+3=(x-8)^2\;\Rightarrow\;x=8+\sqrt{y+3}$ (karena $x>8$).
$C=8,\;D=3$. $10C+D=83$.
\end{pembox}

\begin{soal}{26}
Jika $\cos(2x)=-\dfrac{34}{64}$ dengan $0<x<\pi$,
maka $64\sin(x)$ adalah \ldots
\begin{pilB}
  \item $28$ \item $42$ \item $49$ \item $56$ \item $64$
\end{pilB}
\end{soal}
\begin{pembox}
\jawaban{(D) $64\sin x=56$}
$1-2\sin^2 x=-\tfrac{17}{32}\;\Rightarrow\;\sin^2 x=\tfrac{49}{64}\;\Rightarrow\;\sin x=\tfrac{7}{8}$.
$64\sin x=56$.
\end{pembox}

\begin{soal}{27}
Diketahui $AB$ adalah diameter suatu lingkaran dengan $A=(10,16)$
dan $B=(-2,0)$. Tentukan jari-jari lingkaran dan titik pusat
lingkaran tersebut.
\end{soal}
\begin{pembox}
\jawaban{Jari-jari $r=10$, titik pusat $O=(4,8)$}
Titik pusat: $O=\left(\tfrac{10+(-2)}{2},\tfrac{16+0}{2}\right)=(4,8)$.
$|AB|=\sqrt{12^2+16^2}=20$, sehingga $r=10$.
\end{pembox}

\begin{soal}{28}
Nilai fungsi $f(x)$, $g(x)$, dan turunannya di $x=0$ dan $x=1$
diberikan pada tabel berikut.
\begin{center}
\begin{tabular}{c|cccc}
$x$ & $f(x)$ & $g(x)$ & $f'(x)$ & $g'(x)$ \\
\hline
$0$ & $2$ & $1$ & $3$ & $\tfrac{2}{3}$ \\
$1$ & $3$ & $-5$ & $-\tfrac{1}{4}$ & $-1$ \\
\end{tabular}
\end{center}
Misalkan $h(x)=\dfrac{2f(x)}{g(x)+1}$.
Tentukan kebenaran setiap pernyataan berikut!

\par\smallskip
{\renewcommand{\arraystretch}{1.4}%
\begin{tabular}{@{}p{0.76\linewidth}cc@{}}
\hline
\textbf{Pernyataan} & \textbf{Benar} & \textbf{Salah}\\
\hline
(1)~$h(0)=2$.
  & $\bigcirc$ & $\bigcirc$\\
(2)~Penghitungan $h'(0)$ menggunakan aturan hasil bagi.
  & $\bigcirc$ & $\bigcirc$\\
(3)~$h'(0)=\dfrac{14}{3}$.
  & $\bigcirc$ & $\bigcirc$\\
(4)~$h'(0)=\dfrac{7}{3}$.
  & $\bigcirc$ & $\bigcirc$\\
\hline
\end{tabular}}
\end{soal}
\begin{pembox}
\jawaban{(1)~\textbf{Benar}\quad(2)~\textbf{Benar}\quad(3)~\textbf{Salah}\quad(4)~\textbf{Benar}}
$h(0)=\tfrac{2\cdot2}{1+1}=2$. Aturan hasil bagi:
\[ h'(0)=\frac{2(3)(2)-2(2)(\tfrac{2}{3})}{2^2}=\frac{12-\tfrac{8}{3}}{4}=\frac{\tfrac{28}{3}}{4}=\frac{7}{3}. \]
Jadi $h'(0)=\tfrac{7}{3}$ (bukan $\tfrac{14}{3}$).
\end{pembox}

\begin{soal}{29}
Tabel berikut menyajikan nilai fungsi $f$ dan turunannya di
beberapa titik.
\begin{center}
\begin{tabular}{c|cccc}
$x$ & $0$ & $4$ & $5$ & $23$ \\
\hline
$f(x)$ & $23$ & $0$ & $7$ & $0$ \\
$f'(x)$ & $7$ & $8$ & $3$ & $6$ \\
\end{tabular}
\end{center}
Tentukan nilai
\[
  \lim_{x\to 4}\frac{f(5x+3)}{x-4}.
\]
\end{soal}
\begin{pembox}
\jawaban{$30$}
\konsep L'Hôpital/definisi turunan: saat $x\to4$, $5x+3\to23$ dan $f(23)=0$.

$f(5x+3)\approx f(23)+f'(23)\cdot(5x+3-23)=6\cdot5(x-4)$.
\[ \lim_{x\to4}\frac{f(5x+3)}{x-4}=30. \]
\end{pembox}

\begin{soal}{30}
Perhatikan limit berikut.
\[
  \lim_{x\to 9}\frac{x+5\sqrt{x}-36}{\sqrt{x}-4}
\]
Tentukan kebenaran setiap pernyataan berikut!

\par\smallskip
{\renewcommand{\arraystretch}{1.4}%
\begin{tabular}{@{}p{0.76\linewidth}cc@{}}
\hline
\textbf{Pernyataan} & \textbf{Benar} & \textbf{Salah}\\
\hline
(1)~Saat $x=9$, penyebut $\sqrt{x}-4=-1\neq0$.
  & $\bigcirc$ & $\bigcirc$\\
(2)~Nilai limit dapat dihitung dengan substitusi langsung $x=9$.
  & $\bigcirc$ & $\bigcirc$\\
(3)~Nilai limit tersebut adalah $-12$.
  & $\bigcirc$ & $\bigcirc$\\
(4)~Nilai limit tersebut adalah $12$.
  & $\bigcirc$ & $\bigcirc$\\
\hline
\end{tabular}}
\end{soal}
\begin{pembox}
\jawaban{(1)~\textbf{Benar}\quad(2)~\textbf{Benar}\quad(3)~\textbf{Salah}\quad(4)~\textbf{Benar}}
Penyebut $\sqrt{9}-4=-1\neq0$, sehingga substitusi langsung berlaku:
\[ \frac{9+5(3)-36}{3-4}=\frac{9+15-36}{-1}=\frac{-12}{-1}=12. \]
Nilai limit adalah $+12$, bukan $-12$.
\end{pembox}

\begin{soal}{31}
Jika
\[
  \int_5^8 f(x)\,dx=3
  \qquad\text{dan}\qquad
  \int_8^5 g(x)\,dx=5,
\]
maka nilai
\[
  \int_5^8\bigl[3f(x)-2g(x)\bigr]\,dx
\]
adalah \ldots
\begin{pilB}
  \item $1$ \item $9$ \item $12$ \item $19$ \item $24$
\end{pilB}
\end{soal}
\begin{pembox}
\jawaban{(D) $19$}
Soal: $\int_8^5 g\,dx=5\;\Rightarrow\;\int_5^8 g\,dx=-5$.
\[ 3(3)-2(-5)=9+10=19. \]
\end{pembox}

\begin{soal}{32}
Jika
\[
  x^6-5x^3+6x^2+ax+b=(x^2-1)\,g(x),
\]
maka $10a+b$ adalah \ldots
\begin{pilB}
  \item $35$ \item $39$ \item $43$ \item $47$ \item $51$
\end{pilB}
\end{soal}
\begin{pembox}
\jawaban{(C) $10a+b=43$}
$x^2-1=(x-1)(x+1)$, sehingga $P(1)=0$ dan $P(-1)=0$.
$P(1)=1-5+6+a+b=0\Rightarrow a+b=-2$.
$P(-1)=1+5+6-a+b=0\Rightarrow -a+b=-12$.
Kurangkan: $2a=10\Rightarrow a=5,\;b=-7$.
$10(5)+(-7)=43$.
\end{pembox}

\begin{soal}{33}
Dari data penerimaan mahasiswa ITB didapatkan informasi berikut.
\begin{enumerate}[leftmargin=*,itemsep=2pt,topsep=3pt]
  \item Sebanyak 46\% dari seluruh mahasiswa pernah mengunjungi
        Kampus Ganesha.
  \item Sebanyak 31\% dari seluruh mahasiswa pernah mengunjungi
        Kampus Jatinangor.
  \item Sebanyak 39\% dari seluruh mahasiswa pernah mengunjungi
        Kampus Ganesha tetapi belum pernah mengunjungi Kampus
        Jatinangor.
\end{enumerate}
Jika dipilih satu mahasiswa secara acak, peluang mahasiswa tersebut
\textbf{belum pernah} mengunjungi Kampus Ganesha maupun Kampus
Jatinangor adalah \ldots
\end{soal}
\begin{pembox}
\jawaban{$30\%$}
$P(G)=46\%$, $P(J)=31\%$, $P(G\cap J^c)=39\%$.
$P(G\cap J)=P(G)-P(G\cap J^c)=46\%-39\%=7\%$.
$P(G\cup J)=46\%+31\%-7\%=70\%$.
$P(\text{belum keduanya})=1-70\%=30\%$.
\end{pembox}

\begin{soal}{34}
Misalkan $x=60\sin t$ dan $y=10\cos t$. Jika $t=\dfrac{3\pi}{4}$,
maka $(x+y)\sqrt{2}$ adalah \ldots
\begin{pilB}
  \item $25$ \item $35$ \item $45$ \item $50$ \item $60$
\end{pilB}
\end{soal}
\begin{pembox}
\jawaban{(D) $50$}
$\sin\tfrac{3\pi}{4}=\tfrac{\sqrt{2}}{2}$, $\cos\tfrac{3\pi}{4}=-\tfrac{\sqrt{2}}{2}$.
$x=60\cdot\tfrac{\sqrt{2}}{2}=30\sqrt{2}$;\; $y=10\cdot(-\tfrac{\sqrt{2}}{2})=-5\sqrt{2}$.
$(x+y)\sqrt{2}=25\sqrt{2}\cdot\sqrt{2}=50$.
\end{pembox}

\begin{soal}{35}
Misalkan $P$ adalah bidang yang melalui tiga titik $A(2,8,3)$,
$B(2,0,0)$, dan $C(0,2,0)$. Titik $X(2,2,w)$ berada pada bidang
$P$ jika nilai $w$ adalah \ldots
\end{soal}
\begin{pembox}
\jawaban{$w=\dfrac{3}{4}$}
$\vec{AB}=(0,-8,-3)$, $\vec{AC}=(-2,-6,-3)$. Vektor normal: $\vec{n}=\vec{AB}\times\vec{AC}=(6,6,-16)$.
Persamaan bidang: $6(x-2)+6(y-0)-16(z-0)=0\;\Rightarrow\;6x+6y-16z=12$.
Substitusi $X(2,2,w)$: $12+12-16w=12\;\Rightarrow\;16w=12\;\Rightarrow\;w=\tfrac{3}{4}$.
\end{pembox}

\begin{soal}{36}
Diketahui vektor-vektor $\vec{a}=(1,x,1)$, $\vec{b}=(y,3,-5)$,
dan $\vec{c}=(xy,-12,z)$. Jika vektor $\vec{a}$ tegak lurus
dengan vektor $\vec{b}$ dan vektor $\vec{b}$ sejajar dengan
vektor $\vec{c}$, maka nilai $y+z$ adalah \ldots
\end{soal}
\begin{pembox}
\jawaban{$y+z=37$ (dengan asumsi $y\neq0$)}
$\vec{a}\perp\vec{b}$: $y+3x-5=0$.\quad $\vec{b}\parallel\vec{c}$: $\vec{c}=k\vec{b}$.
Komponen ke-2: $-12=3k\Rightarrow k=-4$.
Komponen ke-3: $z=(-4)(-5)=20$.
Komponen ke-1: $xy=-4y\Rightarrow y(x+4)=0$. Karena $y\neq0$: $x=-4$.
$y=-4(-4)+5-... $ → $y+3(-4)-5=0\Rightarrow y=17$.
$y+z=17+20=37$.
\end{pembox}

\begin{soal}{37}
Jika
\[
  \int_0^3 5f(x)\,dx+\int_0^1 8f(x)\,dx=18
  \qquad\text{dan}\qquad
  \int_0^3 5f(x)\,dx+\int_1^3 8f(x)\,dx=18,
\]
maka tentukan
\[
  \int_1^3 f(x)\,dx.
\]
\end{soal}
\begin{pembox}
\jawaban{$\displaystyle\int_1^3 f(x)\,dx=1$}
Misalkan $A=\int_0^3 f$, $B=\int_0^1 f$, $C=\int_1^3 f$ ($A=B+C$).
Sistem: $5A+8B=18$ dan $5A+8C=18\;\Rightarrow\;B=C$.
$A=2C$. Substitusi: $10C+8C=18\;\Rightarrow\;C=1$.
\end{pembox}

\begin{soal}{38}
Jika $x_1$ dan $x_2$ adalah dua solusi persamaan
\[
  x^{2\log x-24}=\frac{x^{10}}{10^5},
\]
maka $x_1 x_2=10^k$ dengan $k$ adalah \ldots
\begin{pilB}
  \item $5$ \item $7$ \item $12$ \item $17$ \item $19$
\end{pilB}
\end{soal}
\begin{pembox}
\jawaban{(D) $k=17$}
Misalkan $x=10^t$. Maka $x^{2t-24}=x^{10}/10^5$:
$10^{t(2t-24)}=10^{10t-5}\;\Rightarrow\;2t^2-24t=10t-5\;\Rightarrow\;2t^2-34t+5=0$.
Vieta: $t_1+t_2=17$. $x_1x_2=10^{t_1+t_2}=10^{17}$. Jadi $k=17$.
\end{pembox}

\begin{soal}{39}
Misalkan $S_n$ adalah jumlah parsial $n$ suku pertama suatu deret
geometri. Jika $S_2=16$ dan $S_4=160$, maka suku ke-2 deret
tersebut adalah \ldots
\begin{pilB}
  \item $4$ \item $8$ \item $12$ \item $16$ \item $24$
\end{pilB}
\end{soal}
\begin{pembox}
\jawaban{(C) Suku ke-2 $=12$ (untuk $r>0$)}
$S_2=a(1+r)=16$.\; $S_4=a(1+r)(1+r^2)=160$.
$\tfrac{S_4}{S_2}=1+r^2=10\;\Rightarrow\;r^2=9\;\Rightarrow\;r=3$ (positif).
$a(4)=16\;\Rightarrow\;a=4$. Suku ke-2 $=ar=12$.
\end{pembox}

% ===================================================================
%  SOAL BARU  (nomor 40)
% ===================================================================

\begin{soal}{40}
Diketahui fungsi $f(x)=2x^3-3x^2-12x+4$. Nilai maksimum lokal
dari fungsi tersebut adalah \ldots
\begin{pilB}
  \item $-16$ \item $-4$ \item $4$ \item $11$ \item $15$
\end{pilB}
\end{soal}
\begin{pembox}
\jawaban{(D) Nilai maksimum lokal $=11$}
\konsep Titik kritis: $f'(x)=0$; uji tanda atau $f''$ untuk menentukan jenis.

$f'(x)=6x^2-6x-12=6(x-2)(x+1)$.
Titik kritis: $x=-1$ dan $x=2$.
$f''(x)=12x-6$:\; $f''(-1)=-18<0$ (maksimum lokal);\; $f''(2)=18>0$ (minimum lokal).
$f(-1)=2(-1)-3(1)-12(-1)+4=-2-3+12+4=11$.
\end{pembox}

\end{document}
